% !TEX root = ../main.tex
\corrigesChapitre{3}{Réactions d'oxydoréduction}
\corrigesSouples

\begin{blocLivre}{cor-3-1}
\begin{corrige}{1}{Oxydants, réducteurs et couples}
\demarche{Dans un couple Ox/Red, l'oxydant est écrit en premier ; il capte les électrons dans la demi-équation.}
\begin{questions}
\item Oxydants : \res{\ce{Ag+}}, \res{\ce{Fe^{3+}}}, \res{\ce{Al^{3+}}} ; réducteurs : \res{\ce{Ag}}, \res{\ce{Fe^{2+}}}, \res{\ce{Al}}.
\item $\res{\ce{Ag+ + e- = Ag}}$ ; $\res{\ce{Fe^{3+} + e- = Fe^{2+}}}$ ; $\res{\ce{Al^{3+} + 3e- = Al}}$.
\item $\res{\ce{I2 + 2e- = 2I-}}$ ; $\res{\ce{Cl2 + 2e- = 2Cl-}}$.
\item Dans $\ce{Fe^{3+}}/\ce{Fe^{2+}}$, l'ion fer(II) est le \res{\text{réducteur}} ; dans $\ce{Fe^{2+}}/\ce{Fe}$, il est l'\res{\text{oxydant}}.
\end{questions}
\retiens{Dans une demi-équation, les électrons sont toujours du côté de l'oxydant : l'oxydant capte des électrons pour
devenir le réducteur conjugué. Une même espèce peut être l'oxydant d'un couple et le réducteur d'un autre.}
\end{corrige}
\end{blocLivre}
\begin{blocLivre}{cor-3-2}
\begin{corrige}{2}{Des demi-équations en milieu acide}
\demarche{On ajuste l'élément principal, puis O avec $\ce{H2O}$, H avec $\ce{H+}$, et enfin les charges avec des électrons.}
\begin{questions}
\item Étapes : $\ce{Cr2O7^{2-} = 2Cr^{3+}}$ (chrome) ; $+\,7\,\ce{H2O}$ à droite (oxygène) ; $+\,14\,\ce{H+}$ à gauche (hydrogène) ; charges $-2+14=+12$ à
gauche et $+6$ à droite, donc 6 électrons à gauche : $\res{\ce{Cr2O7^{2-} + 14H+ + 6e- = 2Cr^{3+} + 7H2O}}$.
\item $\res{\ce{O2 + 4H+ + 4e- = 2H2O}}$ et $\res{\ce{H2O2 + 2H+ + 2e- = 2H2O}}$.
\item $\res{\ce{NO3- + 4H+ + 3e- = NO + 2H2O}}$.
\item $\res{\ce{O2 + 2H+ + 2e- = H2O2}}$. Le peroxyde d'hydrogène est l'\res{\text{oxydant}} du couple $\ce{H2O2}/\ce{H2O}$ et le
\res{\text{réducteur}} du couple $\ce{O2}/\ce{H2O2}$.
\end{questions}
\piege{Ajuster les charges avant d'avoir ajusté O et H : on obtient un nombre d'électrons faux. L'ordre est toujours le
même : élément principal, puis O avec $\ce{H2O}$, puis H avec $\ce{H+}$, et enfin les charges avec $\mathrm{e}^-$.}
\end{corrige}
\end{blocLivre}
\begin{blocLivre}{cor-3-3}
\begin{corrige}{3}{L'arbre de Diane}
\demarche{Un dépôt métallique et un changement de couleur de la solution signalent un transfert d'électrons.}
\begin{questions}
\item Des \res{\text{cristaux d'argent}} se forment sur le fil, et la solution, incolore au départ, \res{\text{bleuit}}.
\item $\ce{Ag+}/\ce{Ag}$ : $\ce{Ag+ + e- = Ag}$ ; $\ce{Cu^{2+}}/\ce{Cu}$ : $\ce{Cu^{2+} + 2e- = Cu}$.
\item On multiplie la première par 2 : $\res{\ce{Cu + 2Ag+ -> Cu^{2+} + 2Ag}}$. Le \res{\text{cuivre est oxydé}}, les \res{\text{ions argent sont réduits}}.
\item Des ions \res{\ce{Cu^{2+}}}, bleus, passent en solution. Avec quelques gouttes de soude, un prélèvement donne un
\res{\text{précipité bleu}}.
\end{questions}
\schema[corrige]{03-ex03-cor}
\end{corrige}
\end{blocLivre}
\begin{blocLivre}{cor-3-4}
\begin{corrige}{4}{Un métal attaqué par un acide}
\demarche{L'équation donne les proportions : une mole de zinc produit une mole de dihydrogène ; puis $V=n\times V_\mathrm{m}$.}
\begin{questions}
\item On approche une flamme de l'orifice du tube : une \res{\text{petite détonation}} signale le dihydrogène.
\item $\ce{Zn -> Zn^{2+} + 2e-}$ et $\ce{2H+ + 2e- -> H2}$, d'où $\res{\ce{Zn + 2H+ -> Zn^{2+} + H2}}$.
\item $n(\ce{Zn})=\dfrac{0{,}654}{65{,}4}=\res{\qty{1,00e-2}{mol}}$ et $n(\ce{H2})=n(\ce{Zn})=\res{\qty{1,00e-2}{mol}}$.
\item $V=\num{1.00e-2}\times24{,}0=\res{\qty{0,240}{L}}$, soit \qty{240}{mL}.
\end{questions}
\end{corrige}
\end{blocLivre}
\begin{blocLivre}{cor-3-5}
\begin{corrige}{5}{Un clou dans une solution bleue}
\demarche{Les produits identifiés désignent les deux couples : l'un a perdu des électrons, l'autre en a gagné.}
\begin{questions}
\item Le dépôt rouge est du \res{\text{cuivre}} ; le précipité vert révèle les ions \res{\ce{Fe^{2+}}}.
\item Couples \res{\ce{Cu^{2+}}/\ce{Cu}} et \res{\ce{Fe^{2+}}/\ce{Fe}}.
\item $\ce{Fe -> Fe^{2+} + 2e-}$ et $\ce{Cu^{2+} + 2e- -> Cu}$, d'où $\res{\ce{Fe + Cu^{2+} -> Fe^{2+} + Cu}}$.
\item Les ions cuivre(II), bleus, sont consommés, tandis que se forment des ions fer(II), \res{\text{vert pâle}} : la teinte
passe du bleu au vert.
\end{questions}
\end{corrige}
\end{blocLivre}
\begin{blocLivre}{cor-3-6}
\begin{corrige}{6}{Établir des équations}
\demarche{Écrire les deux demi-équations dans le sens de la réaction, égaliser les électrons, additionner.}
\begin{questions}
\item On associe $\ce{MnO4- + 8H+ + 5e- -> Mn^{2+} + 4H2O}$ à cinq fois $\ce{Fe^{2+} -> Fe^{3+} + e-}$ :
\[\res{\ce{MnO4- + 8H+ + 5Fe^{2+} -> Mn^{2+} + 5Fe^{3+} + 4H2O}}.\]
\item On associe $\ce{I2 + 2e- -> 2I-}$ à $\ce{2S2O3^{2-} -> S4O6^{2-} + 2e-}$ :
\[\res{\ce{I2 + 2S2O3^{2-} -> 2I- + S4O6^{2-}}}.\]
\item On associe $\ce{H2O2 + 2H+ + 2e- -> 2H2O}$ à $\ce{2I- -> I2 + 2e-}$ :
\[\res{\ce{H2O2 + 2I- + 2H+ -> I2 + 2H2O}}.\]
\item Charges totales : a) $-1+8+10=+17$ et $2+15=+17$ ; b) $-4$ de chaque côté ; c) $0$ de chaque côté. Les trois
équations sont \res{\text{ajustées}}.
\end{questions}
\retiens{Pour combiner deux demi-équations, on multiplie chacune pour que le nombre d'électrons cédés soit égal au nombre
d'électrons captés ; les électrons disparaissent alors de l'équation. La vérification des charges est un contrôle rapide.}
\end{corrige}
\end{blocLivre}
\begin{blocLivre}{cor-3-7}
\begin{corrige}{7}{L'alcootest chimique}
\demarche{Le dichromate capte 6 électrons, l'éthanol en cède 4 : on multiplie par 2 et par 3 pour en échanger 12.}
\begin{questions}
\item $\res{\ce{Cr2O7^{2-} + 14H+ + 6e- = 2Cr^{3+} + 7H2O}}$.
\item $\res{\ce{CH3COOH + 4H+ + 4e- = CH3CH2OH + H2O}}$.
\item $2\times(\ce{Cr2O7^{2-} + 14H+ + 6e- -> 2Cr^{3+} + 7H2O})$ et $3\times(\ce{CH3CH2OH + H2O -> CH3COOH + 4H+ + 4e-})$. Après
simplification des $\ce{H+}$ et de l'eau :
\[\res{\ce{2Cr2O7^{2-} + 3CH3CH2OH + 16H+ -> 4Cr^{3+} + 3CH3COOH + 11H2O}}.\]
Il y a 4 Cr, 17 O, 34 H et 6 C de chaque côté, et une charge de $+12$ de chaque côté.
\item Les cristaux passent de l'\res{\text{orange au vert}}. La quantité de dichromate réduit est proportionnelle à la
quantité d'éthanol : plus il y a d'alcool, plus la zone verte est longue.
\end{questions}
\end{corrige}
\end{blocLivre}
\begin{blocLivre}{cor-3-8}
\begin{corrige}{8}{Nettoyer l'argenterie}
\demarche{L'aluminium cède 3 électrons, le sulfure d'argent en capte 2 : on échange 6 électrons. Les proportions de l'équation
relient ensuite les quantités.}
\begin{questions}
\item Oxydant : \res{\ce{Ag2S}} ; réducteur : \res{\ce{Al}}.
\item On multiplie $\ce{Ag2S + 2e- -> 2Ag + S^{2-}}$ par 3 et $\ce{Al -> Al^{3+} + 3e-}$ par 2 :
\[\res{\ce{3Ag2S + 2Al -> 6Ag + 2Al^{3+} + 3S^{2-}}}.\]
\item $M(\ce{Ag2S})=2\times107{,}9+32{,}1=\qty{247,9}{g.mol^{-1}}$, $n=\dfrac{0{,}100}{247{,}9}=\qty{4,03e-4}{mol}$ ; $n(\ce{Ag})=2n$ et
$m=2\times\num{4.03e-4}\times107{,}9=\res{\qty{0,0871}{g}}$.
\item D'après l'équation, $n(\ce{Al})=\dfrac{2}{3}\,n(\ce{Ag2S})=\qty{2,69e-4}{mol}$, d'où $m=\res{\qty{7,3}{mg}}$. L'argent du sulfure est \res{\text{redéposé sur
l'objet}} sous forme métallique, alors qu'un polissage arrache le sulfure, et l'argent qu'il contient.
\end{questions}
\piege{Oublier les nombres stœchiométriques dans les calculs de quantités : ici, deux moles d'argent par mole de $\ce{Ag2S}$,
et deux moles d'aluminium pour trois moles de $\ce{Ag2S}$.}
\end{corrige}
\end{blocLivre}
\begin{blocLivre}{cor-3-9}
\begin{corrige}{9}{Vérifier une équation en Python}
\demarche{Un dictionnaire associe une clé à une valeur ; ici, chaque nom d'espèce est associé au couple (composition, charge).}
\begin{questions}
\item \code{especes} associe à chaque espèce sa \res{\text{composition}} (nombre d'atomes de chaque élément) et sa \res{\text{charge}}.
\code{bilan} renvoie le nombre total d'atomes de chaque élément et la charge totale d'un côté de l'équation.
\item La première ligne devient \code{(\{'Mn': 1, 'O': 4, 'H': 4, 'Fe': 5\}, 13)} : il manque 4 atomes H et 4 charges
positives par rapport aux produits. L'équation n'est \res{\text{pas ajustée}}.
\item On modifie les \res{\text{trois dictionnaires}}. Le dictionnaire \code{especes} décrit les quatre espèces, par
exemple avec \code{'Cu2+': (\{'Cu': 1\}, 2)}. Puis \code{reactifs} devient \code{\{'Zn': 1, 'Cu2+': 1\}}, et \code{produits}
devient \code{\{'Zn2+': 1, 'Cu': 1\}}. Les deux lignes affichées sont alors identiques : \code{(\{'Zn': 1, 'Cu': 1\}, 2)}.
\item Une équation peut conserver les éléments sans conserver les charges : $\ce{Fe^{3+} + Cu -> Fe^{2+} + Cu^{2+}}$ porte $+3$ à
gauche et $+4$ à droite. Il faut écrire $\res{\ce{2Fe^{3+} + Cu -> 2Fe^{2+} + Cu^{2+}}}$.
\end{questions}
\end{corrige}
\end{blocLivre}
\begin{blocLivre}{cor-3-10}
\begin{corrige}{10}{La vitamine C, un antioxydant}
\demarche{L'acide ascorbique est un réducteur : il cède 2 électrons ; on l'associe à chaque oxydant.}
\begin{questions}
\item $\res{\ce{C6H6O6 + 2H+ + 2e- = C6H8O6}}$.
\item On associe $\ce{C6H8O6 -> C6H6O6 + 2H+ + 2e-}$ à $\ce{I2 + 2e- -> 2I-}$ :
\[\res{\ce{C6H8O6 + I2 -> C6H6O6 + 2I- + 2H+}}.\]
\item On associe deux fois l'oxydation de l'acide ascorbique à $\ce{O2 + 4H+ + 4e- -> 2H2O}$ ; les ions $\ce{H+}$ se simplifient :
\[\res{\ce{2C6H8O6 + O2 -> 2C6H6O6 + 2H2O}}.\]
\item La vitamine C réagit avec le dioxygène \res{\text{à la place}} des autres constituants de l'aliment, qu'elle protège de
l'oxydation. Dans un jus ouvert, elle est peu à peu \res{\text{consommée par le dioxygène}} de l'air.
\end{questions}
\end{corrige}
\end{blocLivre}
\begin{blocLivre}{cor-3-11}
\begin{corrige}{11}{Ne jamais mélanger eau de Javel et détartrant}
\demarche{Le dichlore est à la fois le réducteur du premier couple et l'oxydant du second : il est le produit des deux
demi-équations.}
\begin{questions}
\item $\res{\ce{2ClO- + 4H+ + 2e- = Cl2 + 2H2O}}$ et $\res{\ce{Cl2 + 2e- = 2Cl-}}$.
\item La réduction $\ce{2ClO- + 4H+ + 2e- -> Cl2 + 2H2O}$ et l'oxydation $\ce{2Cl- -> Cl2 + 2e-}$ s'additionnent en
\[\ce{2ClO- + 2Cl- + 4H+ -> 2Cl2 + 2H2O},\]
qu'on simplifie par 2. Le \res{\text{même élément}}, le chlore, est présent dans l'oxydant et dans le réducteur, qui forment le même produit.
\item $n(\ce{ClO-})=0{,}37\times0{,}100=\qty{0,037}{mol}$, d'où $\res{n(\ce{Cl2})=\qty{0,037}{mol}}$, $V=0{,}037\times24{,}0=\res{\qty{0,89}{L}}$ et
$m=0{,}037\times71{,}0=\res{\qty{2,6}{g}}$.
\item En milieu acide, l'eau de Javel libère un \res{\text{gaz toxique}}, irritant pour les voies respiratoires, même en
petite quantité.
\end{questions}
\end{corrige}
\end{blocLivre}
\begin{blocLivre}{cor-3-12}
\begin{corrige}{12}{Protéger une coque de bateau}
\demarche{Le zinc et le fer cèdent chacun 2 électrons par atome : une mole de zinc oxydée épargne une mole de fer.}
\begin{questions}
\item On associe deux fois $\ce{Fe -> Fe^{2+} + 2e-}$ à $\ce{O2 + 4H+ + 4e- -> 2H2O}$ :
\[\res{\ce{2Fe + O2 + 4H+ -> 2Fe^{2+} + 2H2O}}.\]
\item De même, $\res{\ce{2Zn + O2 + 4H+ -> 2Zn^{2+} + 2H2O}}$.
\item $n(\ce{Zn})=\dfrac{2000}{65{,}4}=\res{\qty{30,6}{mol}}$ ; chaque atome cède 2 électrons : $\res{n(\mathrm{e}^-)=\qty{61,2}{mol}}$.
\item Ces électrons auraient été cédés par $\dfrac{61{,}2}{2}=\qty{30,6}{mol}$ de fer, soit $m=30{,}6\times55{,}8\approx\res{\qty{1,71}{kg}}$ de
fer par an.
\end{questions}
\end{corrige}
\end{blocLivre}
\begin{blocLivre}{cor-3-13}
\begin{corrige}{13}{L'eau oxygénée « à 10 volumes »}
\demarche{Le titre donne $n(\ce{O2})$ par litre ; l'équation de décomposition indique que deux moles de $\ce{H2O2}$ libèrent une
mole de $\ce{O2}$.}
\begin{questions}
\item $\res{\ce{H2O2 + 2H+ + 2e- = 2H2O}}$ et $\res{\ce{O2 + 2H+ + 2e- = H2O2}}$.
\item On additionne la réduction $\ce{H2O2 + 2H+ + 2e- -> 2H2O}$ et l'oxydation $\ce{H2O2 -> O2 + 2H+ + 2e-}$ :
\[\ce{2H2O2 -> 2H2O + O2}.\]
Le peroxyde d'hydrogène est à la fois \res{\text{oxydant et réducteur}} : c'est une \textbf{dismutation}, une même espèce réagissant
sur elle-même.
\item Un litre libère $n(\ce{O2})=\dfrac{10}{22{,}4}=\qty{0,446}{mol}$ ; il contenait donc $2\times0{,}446=\qty{0,893}{mol}$ de $\ce{H2O2}$ :
$\res{c\approx\qty{0,89}{mol.L^{-1}}}$.
\item À \qty{0}{\degreeCelsius}, $10\times0{,}250=\res{\qty{2,5}{L}}$ ; à \qty{20}{\degreeCelsius}, $0{,}446\times0{,}250\times24{,}0=\res{\qty{2,68}{L}}$.
\end{questions}
\end{corrige}
\end{blocLivre}
\begin{blocLivre}{cor-3-14}
\begin{corrige}{14}{Récupérer l'argent d'un bain photographique}
\demarche{Une mole de zinc cède 2 électrons et réduit deux moles d'ions argent : $n(\ce{Zn})=\dfrac{n(\ce{Ag+})}{2}$.}
\begin{questions}
\item $\res{\ce{Zn + 2Ag+ -> Zn^{2+} + 2Ag}}$ ; le zinc est le \res{\text{réducteur}}, il est oxydé.
\item $m=5{,}0\times2{,}0=\res{\qty{10}{g}}$ et $n(\ce{Ag+})=\dfrac{10}{107{,}9}=\res{\qty{0,0927}{mol}}$.
\item $n(\ce{Zn})=\dfrac{0{,}0927}{2}=\qty{0,0463}{mol}$, $m=0{,}0463\times65{,}4=\res{\qty{3,03}{g}}$ au minimum.
\item Un prélèvement du bain traité, additionné de soude, donne un \res{\text{précipité blanc}}.
\end{questions}
\end{corrige}
\end{blocLivre}
\begin{blocLivre}{cor-3-15}
\begin{corrige}{15}{Identifier un métal}
\demarche{Une mole de métal donne une mole de dihydrogène ; la quantité de gaz donne donc celle du métal, puis $M=\dfrac{m}{n}$.}
\begin{questions}
\item $\res{\ce{M + 2H+ -> M^{2+} + H2}}$.
\item $n(\ce{H2})=\dfrac{0{,}096}{24{,}0}=\res{\qty{4,0e-3}{mol}}$.
\item $n(\mathrm{M})=n(\ce{H2})$, donc $M=\dfrac{0{,}262}{\num{4.0e-3}}=\qty{65,5}{g.mol^{-1}}$ : c'est le \res{\text{zinc}}.
\item Le cuivre \res{\text{n'est pas attaqué}} par l'acide chlorhydrique : le dégagement observé l'exclut. Avec la soude, la
solution obtenue doit donner un \res{\text{précipité blanc}}, caractéristique des ions zinc(II).
\end{questions}
\end{corrige}
\end{blocLivre}
